\section{Methodological Design}\label{sec:methodology}

The methodological design of this study is grounded in a separable formulation of the classical spatial gravity framework. Under this formulation, the commuter flow between an origin zone and a destination zone is modeled as a multiplicative product of distinct spatial components. This decomposition allows us to isolate the dominant components of spatial interaction from local, city-specific variables. We structure our investigation around three core hypotheses: (1) aggregate trip-length distributions (TLDs) contain sufficient information to recover city-specific distance-decay parameters without requiring individual origin-destination records; (2) the distance-decay component is the dominant predictive component within the investigated gravity framework; and (3) integrating recovered distance-decay parameters with machine learning models for production and open-data heuristics for attraction preserves strong survey-free target-city flow prediction utility.

\subsection{Gravity Model Foundation}\label{sec:gravity_foundation}

Spatial interaction has long been modeled using gravity-based formulations, which describe travel demand as the combined effect of trip generation, destination attractiveness, and travel impedance. Owing to their physical interpretability and theoretical consistency, gravity models remain one of the most widely adopted frameworks for modeling aggregate human mobility and transportation flows.

Under the singly constrained gravity formulation, the expected flow from origin zone $i$ to destination zone $j$ is expressed as
\begin{equation}\label{eq:decomposition}
T_{ij} = O_i \frac{A_j\, f(d_{ij})}{\sum_{k} A_k\, f(d_{ik})},
\end{equation}
where $O_i$ denotes the total trip production generated by origin zone $i$, $A_j$ represents the attraction of destination zone $j$, $d_{ij}$ is the travel distance between the two zones, and $f(d_{ij})$ is the distance-decay function describing how interaction strength decreases with increasing separation distance. The normalization term ensures that the predicted outflows satisfy the production constraint, $\sum_{j} T_{ij} = O_i$.

A fundamental property of Eq.~(\ref{eq:decomposition}) is its \textbf{separable structure}, in which the three gravity components represent distinct mechanisms governing urban mobility. Origin production captures the total travel demand generated by each zone, destination attraction reflects the spatial distribution of opportunities, and distance decay characterizes travelers' behavioral response to travel cost. Among these components, the distance-decay function is the only term that depends explicitly on inter-zonal distance, whereas origin production and destination attraction are primarily determined by local demographic, socioeconomic, and built-environment characteristics.

This separation has an important implication for survey-free mobility modeling. Rather than reconstructing the entire OD matrix directly, gravity calibration can be reformulated as recovering only the latent behavioral component while estimating production and attraction independently from publicly available spatial information. Consequently, if aggregate mobility observations preserve sufficient information about the distance-dependent interaction mechanism, the gravity model can be calibrated without requiring locally observed OD flows.

\subsection{Aggregate Trip-Length Distributions as a Projection}\label{sec:aggregate_projection}

The gravity formulation presented in the previous section provides a natural interpretation of aggregate trip-length distributions. Rather than representing an independent mobility descriptor, a trip-length distribution can be viewed as the distance-domain projection of the underlying origin--destination interaction process. This interpretation establishes the theoretical connection between aggregate mobility observations and the latent distance-decay mechanism that governs spatial interaction.

Let the travel distance domain be partitioned into a set of distance intervals $\{b\}_{b=1}^{K}$. The aggregate trip-length distribution is defined as
\begin{equation}
P(d_b) = \sum_{i,j} T_{ij}\, \mathbf{1}(d_{ij} \in b),
\end{equation}
where $P(d_b)$ denotes the proportion of trips whose travel distance falls within interval $b$, and $\mathbf{1}(\cdot)$ is the indicator function identifying origin--destination pairs associated with that interval.

Unlike the OD matrix, which explicitly preserves the identities of origins and destinations, the projection removes all location-specific information and retains only the statistical distribution of travel distances. Consequently, aggregate trip-length distributions cannot directly reconstruct the complete OD matrix. Nevertheless, because the gravity formulation separates production, attraction, and distance-dependent behavioral effects, the projection continues to preserve the characteristic statistical signature of the distance-decay component.

Crucially, the observed trip-length distribution is not equivalent to the distance-decay function itself. The projection is simultaneously influenced by urban geometry, the spatial distribution of opportunities, and the unequal number of OD pairs contributing to each distance interval. These structural effects introduce systematic \textbf{exposure biases}, implying that identical behavioral distance-decay functions may produce different aggregate trip-length distributions across cities with different spatial configurations. Therefore, directly fitting a deterrence function to the observed histogram would generally yield biased parameter estimates.

This observation transforms survey-free gravity calibration into an \textbf{inverse inference problem}. Rather than estimating the complete OD matrix, the objective becomes identifying the latent distance-dependent behavioral mechanism whose exposure-corrected projection best explains the observed aggregate trip-length distribution.

\subsection{Distance-Decay Parameterization}\label{sec:decay_function}

To represent travelers' sensitivity to travel distance, we adopt the \textbf{Tanner deterrence function}~\citep{tanner1961,feldman2012alternative}, a flexible two-parameter formulation that combines power-law and exponential decay and has been shown to outperform single-family alternatives across diverse urban settings~\citep{verma2025determinants}. The distance-decay function is expressed as
\begin{equation}\label{eq:tanner_multi}
f(d_{ij};\,\theta) = (d_{ij} + \delta)^{-\alpha} \exp(-\beta d_{ij}),
\end{equation}
where $\theta = (\alpha,\, \beta)$ denotes the unknown behavioral parameters. The power-law exponent $\alpha \ge 0$ controls the concentration of short-to-medium distance trips, whereas the exponential parameter $\beta > 0$ governs the attenuation of long-distance interactions.

To avoid numerical singularities when travel distance approaches zero while maintaining consistent behavior across different spatial partitioning schemes, we introduce an \textbf{adaptive geometric offset}
\begin{equation}\label{eq:delta}
\delta = 0.5\,\bar{R}_{\mathrm{zone}}, \qquad \bar{R}_{\mathrm{zone}} = \frac{1}{N} \sum_{i=1}^{N} \sqrt{\frac{\mathrm{Area}_i}{\pi}}
\end{equation}
where $\bar{R}_{\mathrm{zone}}$ is the average equivalent radius of all spatial zones. Unlike a fixed numerical constant, the adaptive offset automatically scales with the characteristic spatial resolution of the study area, improving numerical stability while preserving the physical interpretation of the recovered distance-decay parameters. During implementation, travel distances are additionally lower-bounded by a small positive constant ($10^{-6}$) to prevent floating-point instability without affecting practical inter-zonal distances.

\subsection{Exposure-Corrected Distance-Decay Recovery}\label{sec:aggregate_calibration}

Building upon the projection formulation above, distance-decay calibration is formulated as a probabilistic inverse inference problem. Instead of directly fitting the observed histogram, the proposed framework explicitly models how aggregate trip-length distributions are generated by the interaction between behavioral deterrence and spatial exposure.

For each distance interval $b$, let $E(b)$ denote the corresponding \textbf{spatial exposure}, defined as the total interaction opportunity associated with all origin--destination pairs whose travel distances belong to interval $b$. Exposure reflects purely structural characteristics of the city---including zone geometry and opportunity distribution---and is independent of travelers' behavioral preferences.

Under the gravity formulation, the expected probability that a trip falls within distance interval $b$ is therefore
\begin{equation}\label{eq:predicted_prob_mle}
\pi_b(\theta) = \frac{E(b)\, f(d_b;\,\theta)}{\sum_{b'} E(b')\, f(d_{b'};\,\theta)},
\end{equation}
which naturally separates behavioral effects from geometric exposure.

Given an observed aggregate trip-length histogram with counts $\{y_b\}$, the observations are modeled as realizations of a multinomial distribution with probabilities defined by Eq.~(\ref{eq:predicted_prob_mle}). The corresponding log-likelihood is
\begin{equation}\label{eq:likelihood_mle}
\log \mathcal{L}(\theta) = \sum_b y_b \log \pi_b(\theta),
\end{equation}
and the city-specific distance-decay parameters are estimated by solving
\begin{equation}
\hat{\theta} = \arg\max_{\theta}\; \log \mathcal{L}(\theta).
\end{equation}
To guarantee the positivity constraints $\alpha \ge 0$, $\beta > 0$, optimization is performed in the unconstrained logarithmic parameter space using the L-BFGS-B algorithm with multiple randomized initializations. The first initialization is obtained from a moment-based estimate of the observed average trip distance, while the remaining runs are generated through Gaussian perturbations to improve robustness against local optima.

Unlike conventional gravity calibration, which requires observed OD matrices, the proposed formulation estimates only the latent behavioral component. By explicitly correcting for spatial exposure, gravity calibration is transformed from a data-intensive flow-fitting problem into a statistically well-defined parameter identification problem requiring only aggregate trip-length observations. The model structure is shown in Fig.~\ref{fig:framework_pipeline}.

\begin{figure*}[t!]
  \centering
  \begin{tikzpicture}[
    font=\footnotesize,
    node distance=0.8cm and 1.1cm,
    block/.style={draw, rectangle, rounded corners, minimum width=2.2cm, minimum height=0.7cm, align=center, thick, fill=blue!5, draw=blue!40},
    cloud/.style={draw, ellipse, minimum width=2.0cm, minimum height=0.7cm, align=center, fill=green!5, draw=green!40},
    agg/.style={draw, ellipse, minimum width=2.2cm, minimum height=0.7cm, align=center, fill=orange!10, draw=orange!60, thick},
    arrow/.style={-latex, thick},
    dasharrow/.style={-latex, thick, dashed, draw=orange!70}
  ]
    % Inputs (open data)
    \node[cloud] (osm) {OSM / Population};
    \node[agg, below=1.5cm of osm] (dist) {Aggregate \\ Trip-Length data};
 
    % Intermediate Models
    \node[block, right=of osm, yshift=0.7cm] (oi) {Outflow GBDT};
    \node[block, right=of osm, yshift=-0.7cm] (aj) {Attraction Heuristic};
    \node[block, right=of dist, xshift=0.3cm] (decay) {Decay MLE \\ $(\hat{\alpha},\,\hat{\beta})$};
 
    % Fusion & Output
    \node[block, right=5cm of osm] (fusion) {Gravity Fusion};
    \node[cloud, below=0.8cm of fusion] (od) {Predicted OD};
 
    % Connections: open-data inputs
    \draw[arrow] (osm.east) -- ++(0.5,0) |- (oi.west);
    \draw[arrow] (osm.east) -- ++(0.5,0) |- (aj.west);
    \draw[arrow] (dist.east) -- (decay.west);
 
    % Merging arrows into Gravity Fusion
    \coordinate (joint) at (6.5, 0);
    \draw[-] (oi.east) -| (joint);
    \draw[-] (aj.east) -| (joint);
    \draw[-] (decay.east) -| (joint);
    \draw[arrow] (joint) -- (fusion.west);
 
    \draw[arrow] (fusion.south) -- (od.north);
 
  \end{tikzpicture}
  \caption{Overview of the aggregate-calibrated gravity model structure. Globally available open-data features (green ellipses) are used to predict trip production ($O_i$) via a GBDT model and to compute destination attraction ($A_j$) via a training-free min-max heuristic. The \textbf{Aggregate Trip-Length Histogram} (orange ellipse, dashed arrow)---which provides aggregate summaries that significantly reduce privacy exposure compared with individual trajectories---feeds the Decay MLE block to recover city-specific parameters $(\hat{\alpha}, \hat{\beta})$ without requiring any individual OD trajectories. The three components are combined via multiplicative gravity fusion to predict survey-free flow matrices ($T_{ij}$).}
  \label{fig:framework_pipeline}
\end{figure*}


\subsection{Transferable Origin Production Estimation (\texorpdfstring{$O_i$}{O\_i})}\label{sec:origin_outflow}

Within the gravity formulation, origin production represents the total number of trips generated by each spatial zone and is conceptually independent of the behavioral interaction mechanism recovered in the previous section. While distance decay governs travelers' destination choices, origin production is primarily determined by relatively stable demographic and built-environment characteristics. This distinction enables the two components to be estimated independently.

For each spatial zone $i$, we construct an open spatial feature vector $\mathbf{x}_i$ consisting of publicly available variables that describe local urban characteristics. These features include population, land-use composition, point-of-interest (POI) distributions, road-network density, and other built-environment indicators derived from open geospatial datasets such as OpenStreetMap and WorldPop. Because these variables are available for virtually all cities without requiring mobility observations, they provide a transferable basis for estimating trip generation.

The mapping between spatial features and origin production is learned from source cities using a Gradient Boosting Decision Tree (GBDT) regression model,
\begin{equation}
\hat{O}_i = g(\mathbf{x}_i),
\end{equation}
where $g(\cdot)$ denotes the learned nonlinear regression function and $\hat{O}_i$ is the predicted trip production for zone $i$. GBDT is adopted because of its strong predictive performance on heterogeneous tabular data, its ability to model nonlinear interactions without extensive feature engineering, and its robustness under cross-city prediction settings. We train the GBDT on a parsimonious set of \textbf{6 core macroscopic features}: total population ($P_i$), total POI count ($\mathrm{POI}_i$), zone area ($\mathrm{Area}_i$), population density ($\rho_{pop, i}$), POI density ($\rho_{poi, i}$), and road density ($\mathrm{rd}_i$). 

These six features were selected from an initial 37-candidate set using a cross-city SHAP stability analysis, showing no performance degradation compared to the full feature set. This parsimonious selection is grounded in the dimensional scaling of urban physics and Occam's razor, thereby preventing the model from overfitting to source-domain geometries. Hyperparameters and details of the SHAP-based feature selection are provided in Appendix~C.

Once trained on source-city data, the regression model is directly applied to the target city without further adaptation or fine-tuning. Consequently, origin production can be estimated under strict zero-shot conditions using only publicly available spatial information, eliminating the need for local travel surveys or historical OD observations.

\subsection{Destination Attraction Estimation (\texorpdfstring{$A_j$}{A\_j})}\label{sec:destination_attraction}

Unlike origin production, destination attraction reflects the relative ability of each spatial zone to attract trips from all origins within the urban system. Its value is therefore determined not only by the intrinsic characteristics of an individual zone but also by the overall spatial distribution of opportunities across the entire city. This dependence makes destination attraction considerably less transferable than origin production and limits the applicability of conventional supervised learning approaches.

Instead of learning attraction from historical destination inflows, the proposed framework estimates attraction directly from publicly available urban opportunity indicators. Let $\mathbf{x}_j$ denote the feature vector associated with destination zone $j$. The attraction score is estimated through a training-free, scale-invariant attraction heuristic:
\begin{equation}\label{eq:attraction_heuristic}
\hat{A}_j = \frac{1}{2} \left[ \mathrm{minmax}(P_j) + \mathrm{minmax}(\mathrm{POI}_j) \right]
\end{equation}
where $\mathrm{minmax}(x_j) = (x_j - \min_k x_k) / (\max_k x_k - \min_k x_k + 10^{-9})$, $P_j$ is the gridded population count, and $\mathrm{POI}_j$ is the OpenStreetMap POI count. This formulation normalizes absolute scale differences across cities while preserving local opportunity hierarchies, ensuring that each indicator contributes proportionally during attraction estimation. 

This analytical estimation strategy offers three practical advantages. First, it is entirely survey-free, requiring only publicly available urban information. Second, it is training-free, avoiding dependence on historical destination-flow observations. Third, because the attraction scores are computed from normalized indicators, the resulting formulation remains scale-invariant, facilitating direct application across cities with substantially different spatial characteristics. In Section~\ref{sec:rq2}, the ablation analysis validates that this training-free heuristic contributes significantly to the downstream prediction accuracy.

\subsection{Survey-Free Origin--Destination Flow Generation}\label{sec:algorithmic_flow}

The preceding sections independently estimate the three components required by the gravity formulation: the city-specific distance-decay function, transferable origin production, and analytical destination attraction. These components are now integrated within the gravity framework to reconstruct the complete origin--destination flow matrix.

Substituting the recovered behavioral parameters and the estimated production and attraction terms into Eq.~(\ref{eq:decomposition}) yields
\begin{equation}\label{eq:predicted_flow_mle}
\hat{T}_{ij} = \hat{O}_i \frac{\hat{A}_j\, f(d_{ij};\,\hat{\theta})}{\sum_{k} \hat{A}_k\, f(d_{ik};\,\hat{\theta})},
\end{equation}
where $\hat{O}_i$ denotes the estimated origin production, $\hat{A}_j$ represents the estimated destination attraction, and $\hat{\theta} = (\hat{\alpha},\, \hat{\beta})$ are the recovered Tanner parameters.

Equation~(\ref{eq:predicted_flow_mle}) preserves the production constraint $\sum_{j} \hat{T}_{ij} = \hat{O}_i$ while allocating trips among competing destinations according to both their relative opportunities and the recovered behavioral response to travel distance. Consequently, each gravity component contributes a distinct aspect of the final prediction. Origin production determines the total travel demand generated by each zone, destination attraction governs the distribution of opportunities across destinations, and the recovered distance-decay function captures travelers' behavioral preferences for spatial interaction.

Unlike conventional gravity calibration, which estimates these components directly from observed OD matrices, PIGF infers each component from the information source most directly associated with its underlying mechanism. Because no target-city OD matrices, household travel surveys, GPS trajectories, or mobile-phone records are required during either calibration or prediction, Eq.~(\ref{eq:predicted_flow_mle}) enables survey-free and zero-shot OD generation using only publicly accessible information.

\subsection{Properties of the Projection--Inference Gravity Framework}

The proposed PIGF possesses several methodological characteristics that distinguish it from existing gravity calibration and supervised OD prediction approaches.

\textbf{Survey-free calibration.} Unlike conventional gravity models that require locally observed OD matrices for parameter calibration, PIGF estimates all gravity components without using any target-city mobility observations. Behavioral information is recovered from aggregate trip-length distributions, while production and attraction are inferred from open spatial datasets.

\textbf{Component-wise inference.} Instead of learning the complete OD matrix as a black-box mapping, PIGF decomposes spatial interaction into three interpretable components---origin production, destination attraction, and distance decay---and estimates each component using the information source most directly associated with its physical or behavioral meaning.

\textbf{Transferability.} All components of PIGF rely exclusively on publicly available or aggregate information, allowing the framework to be directly deployed in previously unseen cities without recalibration using local mobility observations. This property naturally enables zero-shot OD prediction in data-scarce environments.

\begin{table*}[t]
\centering
\tablebodyfont\small
\begin{tabular}{p{0.95\textwidth}}
\toprule
\textbf{Algorithm 1: Projection--Inference Gravity Framework (PIGF)} \\
\midrule
\textbf{Input}: Source cities data $\{\mathbf{X}^{(c)}, \mathbf{T}^{(c)}\}_{c=1}^{C}$; Target city spatial features $\mathbf{X}^{(t)}$; Target city aggregate trip-length histogram $\{y_b\}_{b=1}^{K}$ (or national prior if unavailable) \\
\textbf{Output}: Predicted Target City OD Matrix $\mathbf{\hat{T}}^{(t)}$ \\
\midrule
\textbf{Phase 1: Offline Training (Source Cities)} \\
1. \textbf{Production Model Training}: Train the outflow model $g(\mathbf{x})$ (GBDT) on pooled source-city spatial features $\mathbf{X}^{(c)}$ and ground-truth outflows $O_i^{(c)}$. \\
\midrule
\textbf{Phase 2: Online Calibration \& Zero-Shot Transfer (Target City)} \\
2. \textbf{Decay Calibration}: Recover target-city decay parameters $(\hat{\alpha}, \hat{\beta})$ via MLE on the target city's aggregate distance-bin histogram $\{y_b\}_{b=1}^{K}$. (If unavailable, fall back to the source-derived national prior). \\
3. \textbf{Attraction Mapping}: Compute target city zone attractions using the training-free heuristic: $\hat{A}_j^{(t)} = \frac{1}{2} \left[ \mathrm{minmax}(P_j^{(t)}) + \mathrm{minmax}(\mathrm{POI}_j^{(t)}) \right]$. \\
4. \textbf{Zero-Shot Target Prediction}: \\
\hspace{1em} a. Predict target city outflows: $\hat{O}_i^{(t)} = g(\mathbf{x}_i^{(t)})$. \\
\hspace{1em} b. Compute adaptive target offset: $\delta^{(t)} = 0.5 \times \bar{R}_{\mathrm{zone}}^{(t)}$. \\
\hspace{1em} c. Fuse components and normalize outflows using Eq.~(\ref{eq:predicted_flow_mle}) with recovered parameters $(\hat{\alpha}, \hat{\beta})$. \\
\bottomrule
\end{tabular}
\end{table*}
